\section{Conclusion}
\label{sec:conclusion}

Reading an unfamiliar storage format normally requires a reader written for that format inside every system that wants to consume it. AnyBlox removes this requirement by moving the reader into the data. A decoder travels with the dataset as sandboxed Wasm bytecode, and any engine that implements a single scan interface can decode the format without knowing it in advance. The cost of this design is the sandbox overhead, which the measurements of Gienieczko et al.\ (2025) show to be negligible for the end-to-end workloads tested.

The authors frame the underlying problem as an $N \times M$ problem: $N$ systems each supporting $M$ formats incur $N \times M$ implementation and maintenance effort~\cite{gienieczkoAnyBloxFrameworkSelfDecoding2025}. An abstraction layer between the two sides reduces this to $N + M$, since each engine implements the interface once and each format provides one decoder. The authors draw this analogy explicitly from compiler construction, citing the language-independent intermediate representation of LLVM~\cite{lattnerLLVMCompilationFramework2004}, which decouples $M$ language frontends from $N$ target backends through a shared representation. AnyBlox applies the same decoupling between storage formats and query engines. This mapping is suggestive rather than demonstrated: it explains why the design should work, but the parallel to compilers does not by itself show that formats and engines separate as cleanly as languages and target backends.

This decoupling exemplifies a broader shift toward composable database systems, replacing large, monolithic designs with modular layers. AnyBlox demonstrates this by contributing the storage-format layer and stands alongside efforts modularising other layers: composable-query engines~\cite{lambApacheArrowDataFusion2024}, extensible execution-plan frameworks~\cite{bandleDatabaseTechnologyMasses2021,jungmairDeclarativeSubOperatorsUniversal2023}, portable query optimizers~\cite{alotaibiQueryOptimizerService2024,anneserAutoSteerLearnedQuery2023}, and system-independent persistence~\cite{tsalapatisMemSnapMCheckpointsData2024}. AnyBlox's choice not to define a row-filtering predicate format (see Section~\ref{sec:limitations}) leaves a gap that Substrait could fill. These modular components are complementary, each standardising a specific interface between otherwise independent systems~\cite{gienieczkoAnyBloxFrameworkSelfDecoding2025}.

The authors cite LLVM as a precedent for successful modular frameworks, showing how shared, open representations can transform systems~\cite{gienieczkoAnyBloxFrameworkSelfDecoding2025, lattnerLLVMCompilationFramework2004}. AnyBlox applies this principle to storage formats, and its evaluation demonstrates that a format-agnostic decoding layer can be built with low overhead and integrated into various engines with modest effort. Although the LLVM analogy is plausible, the evidence rests on a single evaluation from the authors. Ultimately, whether modular stacks displace integrated systems requires broader adoption and independent validation. The paper's main result is that the decoding layer is no longer the primary obstacle.